% Generated draft: requires mathematical and source review.
\documentclass[11pt]{article}
\usepackage[margin=25mm]{geometry}
\usepackage{fontspec}
\setmainfont{Latin Modern Roman}
\usepackage{amsmath,amssymb,mathtools,unicode-math}
\setmathfont{Latin Modern Math}
\usepackage{enumitem,needspace,xurl,hyperref}
\setlength{\parindent}{0pt}
\setlength{\parskip}{5pt}
\setlength{\emergencystretch}{3em}
\newcommand{\N}{\mathbb{N}}
\newcommand{\Z}{\mathbb{Z}}
\newcommand{\Q}{\mathbb{Q}}
\newcommand{\R}{\mathbb{R}}
\newcommand{\C}{\mathbb{C}}
\newcommand{\F}{\mathbb{F}}
\newcommand{\A}{\mathbb{A}}
\newcommand{\PP}{\mathbb{P}}
\newcommand{\E}{\mathbb{E}}
\newcommand{\eps}{\varepsilon}
\newcommand{\dd}{\,\mathrm{d}}
\begin{document}
\section*{UnsolvedMath ID: 30002886}
\textbf{Source problem code:} catalogue entry
\subsection{Definitions and mathematical statement}
Let \(X\) be a smooth rigid-analytic space. The supplied record does not specify its ground field, the precise sheaf of differential operators, or the completion used to define co-admissible modules.

For terminology, if a local operator algebra has a Fr\'echet--Stein presentation \(A=\varprojlim_n A_n\), an \(A\)-module \(M\) is co-admissible relative to that presentation when it arises as \(M\cong\varprojlim_n M_n\), where each \(M_n\) is a finitely generated \(A_n\)-module and the transition maps induce isomorphisms
\[
A_n\otimes_{A_{n+1}}M_{n+1}\cong M_n.
\]
This explains the module-theoretic terminology; the source does not identify the local operator algebras or presentations intended on \(X\).

The comparison requested by the source is with algebraic characteristic varieties. For a smooth algebraic variety \(Y\) over a field of characteristic zero and a coherent algebraic \(\mathcal D_Y\)-module \(N\), a good filtration \(F\) gives an associated graded module over \(\operatorname{gr}\mathcal D_Y\cong\operatorname{Sym}_{\mathcal O_Y}\mathcal T_Y\). Its support in the cotangent bundle defines
\[
\operatorname{Char}(N)=\operatorname{Supp}_{T^*Y}(\operatorname{gr}_F N),
\]
independently of the good filtration. This is the comparison model, not an additional hypothesis on the rigid-analytic spaces in the source.

\textbf{Problem.} Develop a satisfactory theory of characteristic varieties for co-admissible \(\mathcal D\)-modules on smooth rigid-analytic spaces, analogous to the characteristic-variety theory for algebraic \(\mathcal D\)-modules.

The source states a research objective rather than a fully specified existence conjecture: it supplies neither the precise rigid-analytic module category nor a target space and axioms for the proposed characteristic varieties. In particular, a cotangent-bundle support construction and particular analogues of algebraic properties cannot be imposed as source-stated requirements without further evidence.
\subsection{Short English statement}
Construct a theory assigning characteristic varieties to co-admissible differential-operator modules on smooth rigid-analytic spaces, comparable to the algebraic theory. The source leaves the ground-field hypotheses, operator completion, target space, and required properties unspecified, so this remains a research-program statement requiring mathematical review.
\subsection{Sources}
\begin{itemize}
\item Local UnsolvedMath registry, \texttt{data/unsolvedmath/problems.json}, permanent ID \(30002886\), source problem code \texttt{catalogue entry}, title ``Characteristic Varieties of Rigid-Analytic D-Modules.'' The supplied record attributes the problem to Oberwolfach workshop discussions, records proposed year 2015, and names no individual proposer.
\item Supplied report citation: \emph{Enveloping Algebras and Geometric Representation Theory} (2016), Oberwolfach Reports, pp.~1385--1447, DOI \texttt{10.4171/owr/2015/25}. \url{https://doi.org/10.4171/owr/2015/25}. The citation's 2016 year and the record's 2015 provenance are retained without reconciliation.
\item Ardakov--Wadsley, \emph{D-modules on rigid analytic spaces I}. \url{https://arxiv.org/abs/1501.02215}.
\item Ardakov--Bode--Wadsley, \emph{D-modules on rigid analytic spaces III}. \url{https://arxiv.org/abs/1904.13280}.
\item \emph{Oberwolfach Report 2015/25}. \url{https://publications.mfo.de/bitstream/handle/mfo/3472/OWR_2015_25.pdf?isAllowed=y&sequence=1}.
\item Literature assessment dated exactly \texttt{2026-08-21}: the supplied curation labels the problem open and states that dimension theory, weak holonomicity, and Kashiwara-type results are available, while a satisfactory general characteristic-variety theory for coadmissible \(\mathcal D\)-modules remains unresolved. This is supplied dated curation derived from the OWR material, not an independently verified literature assessment; its status and cited claims require independent checking.
\end{itemize}
\end{document}
